\documentclass[12pt, a4paper]{article}

\usepackage{xcolor}
\usepackage{geometry}
\usepackage{hyperref}
\usepackage{titlesec}
\usepackage{enumitem}
\usepackage{longtable}
\usepackage{booktabs}
\usepackage{array}
\usepackage{ragged2e}
\usepackage{fancyhdr}
\usepackage{tcolorbox}
\usepackage{amssymb}
\usepackage{amssymb}

\geometry{margin=2.5cm}

\usepackage{xepersian}
\settextfont{Vazirmatn}
\setlatintextfont{Times New Roman}
\setdigitfont{Vazirmatn}

\hypersetup{
    colorlinks=true,
    linkcolor=blue,
    urlcolor=blue
}

\titleformat{\section}{\Large\bfseries\color{blue!60!black}}{\thesection}{1em}{}
\titleformat{\subsection}{\large\bfseries}{\thesubsection}{1em}{}
\titleformat{\subsubsection}{\normalsize\bfseries}{\thesubsubsection}{1em}{}

\setlist{nosep, leftmargin=*}

\pagestyle{fancy}
\fancyhf{}
\fancyhead[R]{\small جزوهٔ مطالعه — جلسهٔ ۵}
\fancyhead[L]{\small Sprint 1 Review}
\fancyfoot[C]{\thepage}
\renewcommand{\headrulewidth}{0.4pt}

\linespread{1.4}

\title{
    \vspace{-1cm}
    \textbf{\Huge جزوهٔ مطالعه}\\[0.3cm]
    \textbf{\LARGE جلسهٔ ۵: Sprint 1 Review}\\[0.5cm]
    \large مهندسی نرم‌افزار
}
\author{
    تهیه‌کننده: ابراهیم فضلی\\
    ترم: ۱۴۰۴
}
\date{\today}

\begin{document}

\maketitle
\thispagestyle{empty}

\vspace{1cm}

\begin{center}
\begin{tabular}{|l|l|}
\hline
\textbf{درس} & مهندسی نرم‌افزار \\
\hline
\textbf{جلسه} & ۵ \\
\hline
\textbf{موضوع} & Sprint 1 Review و Retrospective \\
\hline
\textbf{زمان مطالعه} & ۲-۳ ساعت \\
\hline
\textbf{پیش‌نیاز} & جلسات ۱-۴ \\
\hline
\end{tabular}
\end{center}

\newpage
\tableofcontents
\newpage

% ============================================================
\section{مقدمه}

\subsection{چرا این جلسه مهم است؟}
در جلسات قبل یاد گرفتیم:
\begin{itemize}
    \item \textbf{جلسهٔ ۱:} مهندسی نرم‌افزار چیست
    \item \textbf{جلسهٔ ۲:} چطور نیازمندی‌ها را کشف کنیم
    \item \textbf{جلسهٔ ۳:} فرآیندهای توسعه (Scrum)
    \item \textbf{جلسهٔ ۴:} مدل‌سازی و معماری
\end{itemize}

حالا وقت آن است که ببینیم چه چیزی ساخته‌ایم.

Sprint Review جلسه‌ای است که در آن:
\begin{itemize}
    \item تیم محصول را به مشتری نشان می‌دهد
    \item مشتری بازخورد می‌دهد
    \item دربارهٔ Sprint بعدی تصمیم می‌گیریم
\end{itemize}

\textbf{این جلسه، قلب Scrum است.} بدون آن، Scrum فقط یک سری جلسهٔ بی‌معنی 
است.

\subsection{آنچه در این جزوه یاد می‌گیرید}
\begin{itemize}
    \item Sprint Review چیست و چرا مهم است
    \item تفاوت Sprint Review و Retrospective
    \item چطور Demo مؤثر ارائه دهیم
    \item چطور بازخورد سازنده بگیریم
    \item چطور Retrospective انجام دهیم
    \item چطور Sprint بعدی را برنامه‌ریزی کنیم
\end{itemize}

% ============================================================
\section{Sprint Review چیست؟}

\subsection{تعریف رسمی}
\begin{quote}
\textbf{Sprint Review} یک جلسهٔ غیررسمی در پایان هر Sprint است که در آن تیم 
توسعه محصول را به ذی‌نفعان نشان می‌دهد و بازخورد می‌گیرد.
\end{quote}

\textbf{منبع:} Scrum Guide 2020

\subsection{ویژگی‌های کلیدی}

\begin{center}
\begin{tabular}{|l|p{9cm}|}
\hline
\textbf{ویژگی} & \textbf{توضیح} \\
\hline
مدت & حداکثر ۴ ساعت برای Sprint یک‌ماهه (۱-۲ ساعت برای Sprint دو هفته‌ای) \\
\hline
شرکت‌کنندگان & تیم توسعه + Product Owner + Scrum Master + مشتری + ذی‌نفعان 
\\
\hline
هدف & بازخورد، نه نمره \\
\hline
خروجی & Backlog به‌روز + تصمیم برای Sprint بعدی \\
\hline
فراوانی & پایان هر Sprint \\
\hline
\end{tabular}
\end{center}

\subsection{Sprint Review چه چیزی \textbf{نیست}؟}

\begin{itemize}
    \item ❌ \textbf{گزارش به مدیر نیست} — یک جلسهٔ تعاملی است
    \item ❌ \textbf{ارزیابی عملکرد افراد نیست} — ارزیابی محصول است
    \item ❌ \textbf{ارائهٔ اسلاید نیست} — نمایش محصول کارکننده است
    \item ❌ \textbf{جلسهٔ تصمیم‌گیری نهایی نیست} — بازخورد جمع می‌کنیم
\end{itemize}

\subsection{مثال واقعی}

تصور کنید تیم شما یک فروشگاه آنلاین ساخته است.

\textbf{Sprint Review بد:}
\begin{quote}
«سلام، ما در این Sprint ۸ تا Story را تحویل دادیم. این هم اسلایدهای ما. 
سؤالی هست؟»
\end{quote}

\textbf{Sprint Review خوب:}
\begin{quote}
«سلام. امروز می‌خواهم نشان دهم چطور یک مشتری می‌تواند محصول را پیدا کند و به 
سبد اضافه کند. (Demo زنده) — اینجا صفحهٔ محصولات است. اگر روی فیلتر برند 
کلیک کنیم... (کلیک) — ببینید چطور لیست فیلتر می‌شود. حالا اگر محصول را به 
سبد اضافه کنیم... (کلیک) — سبد بلافاصله به‌روز می‌شود. نظر شما چیست؟ چه چیزی 
را می‌شد بهتر کرد؟»
\end{quote}

% ============================================================
\section{تفاوت Sprint Review و Retrospective}

این دو جلسه اغلب اشتباه گرفته می‌شوند. بگذار تفاوت را واضح کنم:

\begin{center}
\begin{tabular}{|l|l|l|}
\hline
\textbf{ویژگی} & \textbf{Sprint Review} & \textbf{Retrospective} \\
\hline
چه کسی؟ & تیم + مشتری + ذی‌نفعان & فقط تیم \\
\hline
چه چیزی؟ & محصول & فرآیند \\
\hline
هدف & بازخورد روی محصول & بهبود فرآیند \\
\hline
خروجی & Backlog به‌روز & اقدام‌های بهبود \\
\hline
سؤال اصلی & «چه ساختیم؟» & «چطور کار کردیم؟» \\
\hline
زمان & قبل از Retrospective & بعد از Sprint Review \\
\hline
\end{tabular}
\end{center}

\subsection{مثال}

\textbf{Sprint Review:}
\begin{quote}
«مشتری گفت: فیلتر برند خوب است، ولی فیلتر قیمت هم می‌خواهم.»
\textbf{نتیجه:} فیلتر قیمت به Backlog اضافه می‌شود.
\end{quote}

\textbf{Retrospective:}
\begin{quote}
«تیم گفت: ما این Sprint خیلی کد زدیم ولی تست ننوشتیم. باید تست را جدی 
بگیریم.»
\textbf{نتیجه:} اقدام: «از Sprint بعد، برای هر Story حداقل یک تست 
بنویسیم.»
\end{quote}

% ============================================================
\section{Sprint Review در Scrum (جزئیات)}

\subsection{دستور جلسهٔ استاندارد}

\begin{center}
\begin{tabular}{|l|c|p{8cm}|}
\hline
\textbf{مرحله} & \textbf{زمان} & \textbf{توضیح} \\
\hline
۱. خوش‌آمد & ۲ دقیقه & Scrum Master جلسه را باز می‌کند \\
\hline
۲. Sprint Goal & ۳ دقیقه & Product Owner هدف Sprint را یادآوری می‌کند \\
\hline
۳. Demo & ۱۰-۱۵ دقیقه & تیم محصول را نشان می‌دهد \\
\hline
۴. بازخورد & ۵-۱۰ دقیقه & مشتری و ذی‌نفعان نظر می‌دهند \\
\hline
۵. Backlog & ۵ دقیقه & Product Owner Backlog را به‌روز می‌کند \\
\hline
۶. Sprint بعدی & ۵ دقیقه & دربارهٔ Sprint بعدی صحبت می‌شود \\
\hline
\end{tabular}
\end{center}

\subsection{نقش‌ها در Sprint Review}

\subsubsection{Product Owner}
\begin{itemize}
    \item Sprint Goal را یادآوری می‌کند
    \item Backlog را به‌روز می‌کند
    \item تصمیم می‌گیرد چه چیزی در Sprint بعدی باشد
\end{itemize}

\subsubsection{Scrum Master}
\begin{itemize}
    \item جلسه را هماهنگ می‌کند
    \item زمان را کنترل می‌کند
    \item مطمئن می‌شود همه صحبت کنند
\end{itemize}

\subsubsection{تیم توسعه}
\begin{itemize}
    \item Demo را ارائه می‌دهد
    \item به سؤالات فنی جواب می‌دهد
    \item بازخورد را ثبت می‌کند
\end{itemize}

\subsubsection{مشتری و ذی‌نفعان}
\begin{itemize}
    \item Demo را می‌بینند
    \item بازخورد می‌دهند
    \item سؤال می‌پرسند
\end{itemize}

\subsection{قوانین Sprint Review}

\textbf{✅ بایدها:}
\begin{itemize}
    \item محصول \textbf{کارکننده} نشان داده شود
    \item همهٔ ذی‌نفعان دعوت شوند
    \item بازخورد \textbf{مکتوب} ثبت شود
    \item جلسه \textbf{زمان‌بندی} داشته باشد
\end{itemize}

\textbf{❌ نبایدها:}
\begin{itemize}
    \item اسلاید نشان دادن (به‌جای Demo)
    \item بحث فنی طولانی
    \item سرزنش تیم
    \item تصمیم‌گیری بدون مشتری
\end{itemize}

% ============================================================
\section{نحوهٔ ارائهٔ Demo}

\subsection{اصول Demo مؤثر}

\subsubsection{۱. محصول کارکننده نشان بده، نه اسلاید}

\textbf{بد:}
\begin{quote}
«ما این صفحه را ساختیم. این هم اسکرین‌شاتش.»
\end{quote}

\textbf{خوب:}
\begin{quote}
«بگذارید نشان دهم. (کلیک) — این صفحهٔ محصولات است. حالا روی فیلتر کلیک 
می‌کنم...»
\end{quote}

\subsubsection{۲. سناریوی واقعی نشان بده}

\textbf{بد:}
\begin{quote}
«این endpoint لیست محصولات را برمی‌گرداند.»
\end{quote}

\textbf{خوب:}
\begin{quote}
«مریم وارد سایت می‌شود، دنبال یخچال می‌گردد، فیلتر می‌کند، انتخاب می‌کند، به 
سبد اضافه می‌کند.»
\end{quote}

\subsubsection{۳. مسیر اصلی (Happy Path) را نشان بده}
\begin{itemize}
    \item کاربر چه کاری می‌کند؟
    \item چه چیزی می‌بیند؟
    \item چه نتیجه‌ای می‌گیرد؟
\end{itemize}

\subsubsection{۴. محدودیت‌ها را صادقانه بگو}

\textbf{بد:}
\begin{quote}
«همه‌چیز کار می‌کند.»
\end{quote}

\textbf{خوب:}
\begin{quote}
«این قسمت کار می‌کند. این قسمت هنوز در حال تکمیل است. دلیلش این است که...»
\end{quote}

\subsubsection{۵. از قبل تمرین کن}
\begin{itemize}
    \item Demo را ۲-۳ بار تمرین کن
    \item مطمئن شو اینترنت/سرور کار می‌کند
    \item یک نسخهٔ backup (ویدیو) داشته باش
\end{itemize}

\subsection{ساختار Demo (۵ دقیقه)}

\begin{enumerate}
    \item \textbf{معرفی (۳۰ ثانیه):} «امروز نشان می‌دهم چطور مشتری محصول را 
پیدا می‌کند و سفارش می‌دهد.»
    \item \textbf{سناریو (۳ دقیقه):} باز کردن صفحه، جست‌وجو، انتخاب، افزودن 
به سبد، ثبت سفارش
    \item \textbf{نتایج (۱ دقیقه):} «این Storyها تحویل شد: US-01، US-06، 
US-11»
    \item \textbf{محدودیت‌ها (۳۰ ثانیه):} «این قسمت هنوز کامل نیست: ...»
    \item \textbf{دعوت به بازخورد (۳۰ ثانیه):} «نظر شما چیست؟ چه چیزی را 
می‌شد بهتر کرد؟»
\end{enumerate}

\subsection{اشتباهات رایج در Demo}

\begin{center}
\begin{tabular}{|l|l|}
\hline
\textbf{اشتباه} & \textbf{راه‌حل} \\
\hline
اسلاید نشان دادن & محصول واقعی نشان بده \\
\hline
صحبت فنی طولانی & به زبان مشتری صحبت کن \\
\hline
پنهان کردن مشکلات & صادق باش \\
\hline
Demo بدون تمرین & از قبل تمرین کن \\
\hline
Demo طولانی & ۵ دقیقه کافی است \\
\hline
\end{tabular}
\end{center}

% ============================================================
\section{معیارهای ارزیابی Sprint}

\subsection{چه چیزی ارزیابی می‌شود؟}

\begin{center}
\begin{tabular}{|l|c|p{8cm}|}
\hline
\textbf{معیار} & \textbf{وزن} & \textbf{توضیح} \\
\hline
تحویل Storyها & ۴۰٪ & چند تا از Storyها کامل شد \\
\hline
کیفیت کد & ۲۰٪ & تست، Review، خوانایی \\
\hline
مستندات & ۲۰٪ & Vision، Personas، User Stories \\
\hline
ارائه & ۲۰٪ & Demo، وضوح، اعتماد \\
\hline
\end{tabular}
\end{center}

\subsection{چطور نمرهٔ بهتر بگیریم؟}

\subsubsection{۱. Storyها را کامل تحویل بده}
\begin{itemize}
    \item Definition of Done را رعایت کن
    \item تست بنویس
    \item Code Review کن
\end{itemize}

\subsubsection{۲. کیفیت کد را بالا ببر}
\begin{itemize}
    \item کد خوانا بنویس
    \item نام‌گذاری خوب
    \item تست Unit
\end{itemize}

\subsubsection{۳. مستندات را کامل کن}
\begin{itemize}
    \item Vision به‌روز
    \item Personas دقیق
    \item User Stories با Acceptance Criteria
\end{itemize}

\subsubsection{۴. Demo را خوب ارائه بده}
\begin{itemize}
    \item از قبل تمرین کن
    \item محصول کارکننده نشان بده
    \item به بازخورد گوش بده
\end{itemize}

% ============================================================
\section{Definition of Done (با مثال)}

\subsection{تعریف}
\begin{quote}
\textbf{Definition of Done (DoD)} یک توافق تیمی است که مشخص می‌کند یک Story 
چه زمانی «تمام» است.
\end{quote}

\subsection{DoD نمونه}

یک Story وقتی «تمام» است که:
\begin{itemize}
    \item[$\square$] کد نوشته شده و Merge شده
    \item[$\square$] تست Unit دارد و پاس می‌شود
    \item[$\square$] Code Review شده
    \item[$\square$] مستندات به‌روز است
    \item[$\square$] در محیط Demo کار می‌کند
    \item[$\square$] Acceptance Criteria برآورده شده
\end{itemize}

\subsection{مثال واقعی}

\textbf{Story:} US-01 — نمایش محصولات

\textbf{Acceptance Criteria:}
\begin{itemize}
    \item[$\square$] صفحهٔ محصولات با ۲۰ آیتم نمایش داده شود
    \item[$\square$] هر آیتم شامل تصویر، نام، قیمت باشد
    \item[$\square$] بارگذاری < ۲ ثانیه باشد
\end{itemize}

\textbf{Definition of Done:}
\begin{itemize}
    \item[$\square$] کد در \lr{feature/products} نوشته شد
    \item[$\square$] تست \lr{products.test.js} نوشته شد
    \item[$\square$] PR باز شد و Review شد
    \item[$\square$] مستندات \lr{user-stories.md} به‌روز شد
    \item[$\square$] در \lr{localhost:3000} کار می‌کند
\end{itemize}

\textbf{نتیجه:} Story «تمام» است. ✅

% ============================================================
\section{Retrospective (با مثال)}

\subsection{تعریف}
\begin{quote}
\textbf{Retrospective} جلسه‌ای داخلی برای تیم است که در آن دربارهٔ فرآیند 
Sprint قبل صحبت می‌کنند و اقدام‌های بهبود تعیین می‌کنند.
\end{quote}

\subsection{ویژگی‌ها}

\begin{center}
\begin{tabular}{|l|p{9cm}|}
\hline
\textbf{ویژگی} & \textbf{توضیح} \\
\hline
مدت & ۱-۱.۵ ساعت \\
\hline
شرکت‌کنندگان & فقط تیم \\
\hline
هدف & بهبود فرآیند \\
\hline
خروجی & اقدام‌های قابل اجرا \\
\hline
فضا & امن، بدون سرزنش \\
\hline
\end{tabular}
\end{center}

\subsection{قوانین Retrospective}

\textbf{✅ بایدها:}
\begin{itemize}
    \item صادقانه صحبت کن
    \item به دیگران گوش بده
    \item روی فرآیند تمرکز کن، نه افراد
    \item اقدام‌های \textbf{قابل اندازه‌گیری} تعیین کن
\end{itemize}

\textbf{❌ نبایدها:}
\begin{itemize}
    \item سرزنش کردن
    \item بحث فنی طولانی
    \item اقدام‌های مبهم («بهتر کار کنیم»)
    \item جلسهٔ شکایت
\end{itemize}

\subsection{مثال واقعی}

\textbf{تیم هوم‌استور بعد از Sprint 1:}

\textbf{Start (شروع کنیم):}
\begin{itemize}
    \item نوشتن تست از ابتدا (TDD)
    \item جلسهٔ Daily منظم
    \item استفاده از GitHub Issues
\end{itemize}

\textbf{Stop (متوقف کنیم):}
\begin{itemize}
    \item کد زدن بدون تست
    \item Commitهای بزرگ
    \item کار انفرادی
\end{itemize}

\textbf{Continue (ادامه دهیم):}
\begin{itemize}
    \item Code Review
    \item مستندسازی
    \item استفاده از AI با شفافیت
\end{itemize}

\textbf{نتیجه:} ۳ اقدام مشخص برای Sprint 2.

% ============================================================
\section{سؤالات کلیدی برای مطالعه}

\subsection{سطح ۱ (پایه)}

\subsubsection{سؤال ۱: Sprint Review چیست؟ در چه زمانی برگزار می‌شود؟}
\textbf{پاسخ:} Sprint Review یک جلسهٔ غیررسمی در پایان هر Sprint است که در 
آن تیم محصول را به ذی‌نفعان نشان می‌دهد و بازخورد می‌گیرد. در پایان هر Sprint 
برگزار می‌شود (معمولاً هر ۲-۴ هفته).

\subsubsection{سؤال ۲: تفاوت Sprint Review و Retrospective چیست؟}
\textbf{پاسخ:}
\begin{itemize}
    \item \textbf{Sprint Review:} با مشتری، دربارهٔ محصول، بازخورد محصول
    \item \textbf{Retrospective:} فقط تیم، دربارهٔ فرآیند، بهبود فرآیند
\end{itemize}

\subsubsection{سؤال ۳: Definition of Done چیست؟}
\textbf{پاسخ:} یک توافق تیمی است که مشخص می‌کند یک Story چه زمانی «تمام» 
است. شامل: کد نوشته شده، تست دارد، Code Review شده، مستندات به‌روز، در Demo 
کار می‌کند.

\subsubsection{سؤال ۴: چه کسانی در Sprint Review شرکت می‌کنند؟}
\textbf{پاسخ:} تیم توسعه، Product Owner، Scrum Master، مشتری، و ذی‌نفعان.

\subsubsection{سؤال ۵: چرا Sprint Review مهم است؟}
\textbf{پاسخ:}
\begin{itemize}
    \item مطمئن می‌شویم در مسیر درست هستیم
    \item مشتری زودتر بازخورد می‌دهد
    \item مشکلات زودتر کشف می‌شوند
    \item تیم انگیزه می‌گیرد
    \item شفافیت ایجاد می‌شود
\end{itemize}

\subsection{سطح ۲ (متوسط)}

\subsubsection{سؤال ۶: چطور یک Demo مؤثر ارائه دهیم؟}
\textbf{پاسخ:}
\begin{enumerate}
    \item محصول کارکننده نشان بده، نه اسلاید
    \item سناریوی واقعی نشان بده
    \item مسیر اصلی (Happy Path) را نشان بده
    \item محدودیت‌ها را صادقانه بگو
    \item از قبل تمرین کن
\end{enumerate}

\subsubsection{سؤال ۷: چرا اسلاید نشان دادن در Sprint Review بد است؟}
\textbf{پاسخ:} چون Sprint Review برای \textbf{دیدن محصول کارکننده} است، نه 
شنیدن گزارش. اسلاید نشان می‌دهد که تیم اعتماد ندارد محصول را نشان دهد، یا 
محصول آماده نیست. همچنین مشتری نمی‌تواند بازخورد واقعی بدهد.

\subsubsection{سؤال ۸: Start/Stop/Continue چیست؟ چطور استفاده کنیم؟}
\textbf{پاسخ:} یک روش ساده برای Retrospective:
\begin{itemize}
    \item \textbf{Start:} چه کاری را شروع کنیم؟
    \item \textbf{Stop:} چه کاری را متوقف کنیم؟
    \item \textbf{Continue:} چه کاری را ادامه دهیم؟
\end{itemize}
هر بخش حداقل ۱-۲ مورد. اقدام‌ها باید قابل اندازه‌گیری باشند.

\subsubsection{سؤال ۹: چطور بازخورد سازنده بگیریم؟}
\textbf{پاسخ:}
\begin{itemize}
    \item سؤال باز بپرس: «چه چیزی را می‌شد بهتر کرد؟»
    \item به جای دفاع، گوش بده
    \item بازخورد را مکتوب کن
    \item از بازخورد تشکر کن
    \item اقدام مشخص تعیین کن
\end{itemize}

\subsubsection{سؤال ۱۰: چطور اقدام‌های Retrospective را پیگیری کنیم؟}
\textbf{پاسخ:}
\begin{itemize}
    \item هر اقدام یک \textbf{مالک} داشته باشد
    \item هر اقدام \textbf{معیار موفقیت} داشته باشد
    \item در Sprint بعد، بررسی شود
    \item اگر کار نکرد، دلیلش پیدا شود
    \item اگر کار کرد، جشن گرفته شود
\end{itemize}

\subsection{سطح ۳ (پیشرفته)}

\subsubsection{سؤال ۱۱: اگر مشتری در Sprint Review گفت «همه‌چیز خوب است» 
ولی بعداً ناراضی بود، چه کار می‌کنید؟}
\textbf{پاسخ:}
\begin{itemize}
    \item سؤال‌های عمیق‌تر بپرسید: «اگر بخواهید یک چیزی را عوض کنید، چه؟»
    \item سناریوهای واقعی‌تر نشان دهید
    \item از مشتری بخواهید خودش محصول را امتحان کند
    \item بازخورد غیرمستقیم بگیرید (نظرسنجی)
    \item در Retrospective بررسی کنید چرا بازخورد نگرفتیم
\end{itemize}

\subsubsection{سؤال ۱۲: چطور مطمئن شویم Sprint Review به یک جلسهٔ شکایت 
تبدیل نمی‌شود؟}
\textbf{پاسخ:}
\begin{itemize}
    \item قوانین جلسه را از اول مشخص کنید
    \item Scrum Master زمان و فضا را مدیریت کند
    \item روی \textbf{محصول} تمرکز کنید، نه افراد
    \item بازخورد را \textbf{مکتوب} کنید
    \item هر انتقاد را به یک \textbf{اقدام} تبدیل کنید
\end{itemize}

\subsubsection{سؤال ۱۳: اگر تیم در Sprint Review شکست خورد، چطور برخورد 
کنیم؟}
\textbf{پاسخ:}
\begin{itemize}
    \item \textbf{سرزنش نکنید} — شکست فرصت یادگیری است
    \item در Retrospective صادقانه بررسی کنید
    \item دلایل واقعی را پیدا کنید (نه بهانه)
    \item اقدام‌های مشخص تعیین کنید
    \item Sprint بعدی را واقع‌بینانه‌تر برنامه‌ریزی کنید
\end{itemize}

\subsubsection{سؤال ۱۴: چطور Sprint Review را برای تیم‌های دور از هم 
(ریموت) برگزار کنیم؟}
\textbf{پاسخ:}
\begin{itemize}
    \item \textbf{ابزار:} Zoom / Google Meet + Miro
    \item \textbf{Demo:} اشتراک صفحه + کنترل از راه دور
    \item \textbf{بازخورد:} Miro board برای نوشتن همزمان
    \item \textbf{ضبط:} جلسه را ضبط کنید برای غایبان
    \item \textbf{Async:} اگر زمان هماهنگ سخت است، ویدیوی Demo بفرستید
\end{itemize}

% ============================================================
\section{تمرین‌ها}

\subsection{تمرین ۱: تحلیل Sprint Review}
یک ویدیو از Sprint Review واقعی پیدا کنید و تحلیل کنید:
\begin{itemize}
    \item چه چیزی خوب بود؟
    \item چه چیزی بد بود؟
    \item چه چیزی یاد گرفتید؟
\end{itemize}

\subsection{تمرین ۲: نوشتن DoD}
برای پروژهٔ خودتان، Definition of Done بنویسید:
\begin{itemize}
    \item حداقل ۵ مورد
    \item قابل اندازه‌گیری
    \item برای همهٔ Storyها
\end{itemize}

\subsection{تمرین ۳: Retrospective}
یک Retrospective برای Sprint 1 خودتان انجام دهید:
\begin{itemize}
    \item Start / Stop / Continue
    \item حداقل ۳ اقدام
    \item مالک مشخص
\end{itemize}

\subsection{تمرین ۴: Demo}
یک Demo ۵ دقیقه‌ای برای پروژهٔ خودتان آماده کنید:
\begin{itemize}
    \item سناریوی واقعی
    \item محصول کارکننده
    \item دعوت به بازخورد
\end{itemize}

% ============================================================
\section{منابع}

\subsection{منابع اصلی}

\begin{center}
\begin{tabular}{|l|l|c|}
\hline
\textbf{منبع} & \textbf{فصل} & \textbf{صفحه} \\
\hline
Scrum Guide 2020 & Sprint Review & ۱۰-۱۲ \\
\hline
Sommerville — Software Engineering & فصل ۲ & ۳۵-۴۵ \\
\hline
\end{tabular}
\end{center}

\subsection{منابع تکمیلی}
\begin{itemize}
    \item \textbf{ویدیو:} Sprint Review in 5 minutes
    \item \textbf{مقاله:} How to run effective Sprint Review
    \item \textbf{کتاب:} Agile Retrospectives
\end{itemize}

\subsection{ابزارها}
\begin{itemize}
    \item \textbf{Miro} — برای Retrospective آنلاین
    \item \textbf{EasyRetro} — رایگان و ساده
    \item \textbf{FunRetro} — رایگان
\end{itemize}

% ============================================================
\section{جمع‌بندی}

\subsection{سه نکتهٔ کلیدی}

\begin{enumerate}
    \item \textbf{Sprint Review = نمایش، نه گزارش}
    \begin{itemize}
        \item محصول کارکننده نشان بده
        \item بازخورد بگیر
        \item Backlog را به‌روز کن
    \end{itemize}
    \item \textbf{Retrospective = بهبود فرآیند، نه ارزیابی افراد}
    \begin{itemize}
        \item صادقانه صحبت کن
        \item اقدام‌های قابل اندازه‌گیری
        \item پیگیری کن
    \end{itemize}
    \item \textbf{بازخورد = بهترین دوست تیم}
    \begin{itemize}
        \item زودتر بازخورد بگیر
        \item صادقانه گوش بده
        \item اقدام کن
    \end{itemize}
\end{enumerate}

\subsection{سؤال نهایی}

\begin{quote}
\textbf{اگر امروز Sprint Review داشتید، چه چیزی نشان می‌دادید؟}
\end{quote}

\vspace{1cm}

\begin{center}
\textbf{\Large پایان جزوه}\\[0.5cm]
\textbf{\Large موفق باشید!}
\end{center}

\end{document}
